\paragraph{Limitations}
A key limitation of our study lies in the use of a LLM as the verifier for the correctness of reasoning CoTs, due to the prohibitive cost of manually checking a large volume of generated reasoning paths. Moreover, our theorem only explains the optimization process of RLVR but provides no guarantee for its generalization. We merely observe the generalization empirically. 
Due to space limitations, further discussions on the implications of our findings are deferred to Appendix~\ref{supp_sec:discuss}.

\section{Conclusion}
\label{sec:conclu}

In this work, we address the fundamental problem of whether RLVR genuinely incentivizes novel reasoning in base LLMs. Through empirical evaluations and theoretical analysis, we justify a new perspective: RLVR implicitly incentivizes correct reasoning. Moreover, our analyses on training dynamics and CoT quality further confirm that the reasoning CoTs after RLVR are fundamentally different and can even help to replicate similar capabilities simply via supervised learning.

We hope these findings can not only resolve conflicting conclusions in prior work but also illuminate the untapped potential of RLVR in aligning LLMs with human reasoning systems. We envision a promising future where RLVR serves as a cornerstone for developing LLMs that learn through interaction, self-correction, and verifiable reasoning.
